\textbf{Scale.} One PDE ($\nu=0.01$), one data distribution, 400 training pairs, 100 test pairs per seed, and five seeds (three for sweeps). Welch tests on five values are fragile and the very large effect sizes are a consequence of the huge gaps, not of precise estimation. The seed changes the data, the initialisation and the batch order together, so we cannot separate their contributions.

\textbf{The zero-shot test is easy.} Initial conditions contain only 63 Fourier modes and every grid resolves the solution, so a spectral model sees the same function at all resolutions. Real super-resolution (content above the training Nyquist frequency, shocks at smaller $\nu$) is not tested, and our results do not support or contradict the aliasing concerns raised in the literature~\cite{bartolucci2023representation}. Conclusions about MLP and DeepONet depend on our stride-subsampling rule; other conventions (e.g.\ interpolating onto the sensors) would change their numbers.

\textbf{Baselines.} All models are reimplemented from the published descriptions, with one tuned hyperparameter (learning rate) and a common 100-epoch budget. DeepONet and the CNN are under-fit at that budget; we ran longer training only for DeepONet and the MLP, with three seeds, and not for the CNN, whose 49-point receptive field is a second, untested candidate cause of its error. Our DeepONet does not reproduce the near-parity with the FNO reported for simple settings~\cite{lu2021comprehensive}: it is the plain architecture without the extensions used there, it overfits when trained longer (Table~\ref{tab:abl}), and no published error value was used as a reference, so its numbers are a lower bound on DeepONet quality. The training-set-size sweep confounds data size with the number of gradient steps. Larger widths, learning-rate schedules, normalisation layers, or a U-Net or CNO~\cite{raoni2023convolutional} baseline might narrow the gap. Training times are wall-clock on a shared machine and are only indicative. We did not search the FNO's depth, width or optimiser beyond the learning rate. The learning-rate grid was extended once after seeing the FNO's tuning results on the validation split; this affects the FNO only through validation data. The extension was not applied symmetrically at the time: the CNN's point at the new rate was run only in the revision (Table~\ref{tab:tune}), and the FNO's selected rate still sits at the edge of the extended grid, so the FNO may be slightly under-tuned as well. Each tuning cell is a single run on one seed.

\textbf{What a full study would add.} More PDEs and viscosities, out-of-distribution initial conditions, error beyond relative $L_2$ (e.g.\ spectra, shock location), more seeds and matched-convergence training, and the other neural operators discussed above.
